% 各章のファイルは，main.texの\includeで読み込む．
% 章を追加したら，main.texに\include{contents/ファイル名}を追加する．
% ラベルは「環境名:識別子」の形にする(例：Def:group，Thm:lagrange，sec:group)．
\chapter{章の題名} \label{chap:sample}

\section{節の題名} \label{sec:sample}

\begin{Def} \label[Def]{Def:group}
	集合\(G\)と写像\(\cdot \colon G \times G \to G\)の組\(\paren{G, \cdot}\)が%
	\term[ぐん]{群}であるとは，以下の条件をすべて満たすことをいう：
	\begin{enumerate}
		\item 任意の\(a, b, c \in G\)に対し\(\paren{a \cdot b} \cdot c = a \cdot \paren{b \cdot c}\)である． \label{item:group-assoc}
		\item ある\(e \in G\)が存在し，任意の\(a \in G\)に対し\(a \cdot e = e \cdot a = a\)である．
		\item 任意の\(a \in G\)に対し，ある\(b \in G\)が存在して\(a \cdot b = b \cdot a = e\)である．
	\end{enumerate}
\end{Def}

\begin{Ex} \label[Ex]{Ex:integer-group}
	\(\paren{\Z, +}\)は群である．
\end{Ex}

\begin{Thm} \label[Thm]{Thm:unit-unique}
	群の単位元は一意に定まる．
\end{Thm}

\begin{proof}
	\(e, e'\)がともに単位元であるとすると，\(e = e \cdot e' = e'\)である．
\end{proof}

\begin{Note}
	\Cref{Def:group}の条件\cref{item:group-assoc}を結合法則という．
	\Cref{Thm:unit-unique}の証明では，結合法則を使っていない．
\end{Note}

群論の標準的な教科書として\cite{sample2026}を挙げる．

\begin{Que} \label[Que]{Que:inverse-unique}
	群の各元の逆元は一意に定まることを示せ．
\end{Que}
